% Options for packages loaded elsewhere
\PassOptionsToPackage{unicode}{hyperref}
\PassOptionsToPackage{hyphens}{url}
\documentclass[
]{article}
\usepackage{xcolor}
\usepackage{amsmath,amssymb}
\setcounter{secnumdepth}{-\maxdimen} % remove section numbering
\usepackage{iftex}
\ifPDFTeX
  \usepackage[T1]{fontenc}
  \usepackage[utf8]{inputenc}
  \usepackage{textcomp} % provide euro and other symbols
\else % if luatex or xetex
  \usepackage{unicode-math} % this also loads fontspec
  \defaultfontfeatures{Scale=MatchLowercase}
  \defaultfontfeatures[\rmfamily]{Ligatures=TeX,Scale=1}
\fi
\usepackage{lmodern}
\ifPDFTeX\else
  % xetex/luatex font selection
\fi
% Use upquote if available, for straight quotes in verbatim environments
\IfFileExists{upquote.sty}{\usepackage{upquote}}{}
\IfFileExists{microtype.sty}{% use microtype if available
  \usepackage[]{microtype}
  \UseMicrotypeSet[protrusion]{basicmath} % disable protrusion for tt fonts
}{}
\makeatletter
\@ifundefined{KOMAClassName}{% if non-KOMA class
  \IfFileExists{parskip.sty}{%
    \usepackage{parskip}
  }{% else
    \setlength{\parindent}{0pt}
    \setlength{\parskip}{6pt plus 2pt minus 1pt}}
}{% if KOMA class
  \KOMAoptions{parskip=half}}
\makeatother
\usepackage{longtable,booktabs,array}
\newcounter{none} % for unnumbered tables
\usepackage{calc} % for calculating minipage widths
% Correct order of tables after \paragraph or \subparagraph
\usepackage{etoolbox}
\makeatletter
\patchcmd\longtable{\par}{\if@noskipsec\mbox{}\fi\par}{}{}
\makeatother
% Allow footnotes in longtable head/foot
\IfFileExists{footnotehyper.sty}{\usepackage{footnotehyper}}{\usepackage{footnote}}
\makesavenoteenv{longtable}
\usepackage{graphicx}
\makeatletter
\newsavebox\pandoc@box
\newcommand*\pandocbounded[1]{% scales image to fit in text height/width
  \sbox\pandoc@box{#1}%
  \Gscale@div\@tempa{\textheight}{\dimexpr\ht\pandoc@box+\dp\pandoc@box\relax}%
  \Gscale@div\@tempb{\linewidth}{\wd\pandoc@box}%
  \ifdim\@tempb\p@<\@tempa\p@\let\@tempa\@tempb\fi% select the smaller of both
  \ifdim\@tempa\p@<\p@\scalebox{\@tempa}{\usebox\pandoc@box}%
  \else\usebox{\pandoc@box}%
  \fi%
}
% Set default figure placement to htbp
\def\fps@figure{htbp}
\makeatother
\ifLuaTeX
  \usepackage{luacolor}
  \usepackage[soul]{lua-ul}
\else
  \usepackage{soul}
\fi
\setlength{\emergencystretch}{3em} % prevent overfull lines
\providecommand{\tightlist}{%
  \setlength{\itemsep}{0pt}\setlength{\parskip}{0pt}}
\usepackage{bookmark}
\IfFileExists{xurl.sty}{\usepackage{xurl}}{} % add URL line breaks if available
\urlstyle{same}
\hypersetup{
  hidelinks,
  pdfcreator={LaTeX via pandoc}}

\author{}
\date{}

\begin{document}

\textbf{Java Raytracer}

Report

\textbf{Multimedia \& Computer Graphics}

Universidad Panamericana \textbar{} 2026

Santiago Macias Fernandez

\section{}\label{section}

\section{}\label{section-1}

\section{}\label{section-2}

\section{}\label{section-3}

\section{}\label{section-4}

\section{}\label{section-5}

\section{}\label{section-6}

\section{}\label{section-7}

\section{\texorpdfstring{\textbf{Introduction}}{Introduction}}\label{introduction}

This project is basically a raytracer built from scratch in Java. A
raytracer is a rendering technique that simulates how light actually
works in the real world, instead of rasterizing polygons like a regular
game engine does, it shoots rays from the camera into the scene and
calculates exactly what each ray hits and how it bounces around.

The result is that you get physically accurate shadows, reflections, and
refractions kind of for free, because you\textquotesingle re actually
simulating the physics of light. The downside is that
it\textquotesingle s way slower than rasterization, but for images it
gives much better results.

This report walks through how the project evolved from version 0.1 to
version 1.0, what was added in each stage, and what the important
technical decisions were. The codebase ended up being about 20 Java
files and can render complex 3D scenes with textures, normal maps,
reflections, and refractions.

\section{\texorpdfstring{\textbf{v0.1 ---
Spheres}}{v0.1 --- Spheres}}\label{v0.1-spheres}

{\def\LTcaptype{none} % do not increment counter
\begin{longtable}[]{@{}
  >{\raggedright\arraybackslash}p{(\linewidth - 2\tabcolsep) * \real{0.1923}}
  >{\raggedright\arraybackslash}p{(\linewidth - 2\tabcolsep) * \real{0.8077}}@{}}
\toprule\noalign{}
\begin{minipage}[b]{\linewidth}\raggedright
\textbf{v0.1}
\end{minipage} & \begin{minipage}[b]{\linewidth}\raggedright
\textbf{Ray-Sphere Intersection}
\end{minipage} \\
\begin{minipage}[b]{\linewidth}\raggedright
\textbf{Files changed}
\end{minipage} & \begin{minipage}[b]{\linewidth}\raggedright
Raytracer.java, Sphere.java, Vector3D.java, Ray.java, Camera.java,
Scene.java, Object3D.java, Intersection.java
\end{minipage} \\
\midrule\noalign{}
\endhead
\bottomrule\noalign{}
\endlastfoot
\end{longtable}
}

The very first version was just getting a single colored pixel on
screen. Before anything else, the project needed a solid math
foundation, which meant implementing Vector3D.java with all the
operations that a raytracer needs: add, subtract, scalar multiply, dot
product, cross product, magnitude, and normalize.

\subsubsection{\texorpdfstring{\textbf{How ray casting
works}}{How ray casting works}}\label{how-ray-casting-works}

The core idea is simple: for every pixel on screen, shoot a ray from the
camera position through that pixel into the 3D scene. If the ray hits an
object, color the pixel. If it doesn\textquotesingle t, color it black
(background).

The Camera class converts a pixel coordinate (x, y) into a normalized
ray direction. It uses the field-of-view angle and the aspect ratio to
map screen pixels to directions in 3D space. The formula is:

{\def\LTcaptype{none} % do not increment counter
\begin{longtable}[]{@{}
  >{\raggedright\arraybackslash}p{(\linewidth - 0\tabcolsep) * \real{1.0000}}@{}}
\toprule\noalign{}
\begin{minipage}[b]{\linewidth}\raggedright
px = (2 * (x + 0.5) / width - 1) * aspectRatio * tan(fov/2)

py = (1 - 2 * (y + 0.5) / height) * tan(fov/2)

direction = normalize(px, py, -1) // -1 because camera looks towards -Z
\end{minipage} \\
\midrule\noalign{}
\endhead
\bottomrule\noalign{}
\endlastfoot
\end{longtable}
}

For spheres, the intersection math is a quadratic equation. The ray is
parameterized as P(t) = origin + t*direction, and you substitute that
into the sphere equation (\textbar\textbar P -
center\textbar\textbar\^{}2 = radius\^{}2) and solve for t. If t is
positive, the sphere is in front of the camera. The smaller t is the
closer intersection (front face of the sphere).

\includegraphics[width=4.25521in,height=3.19141in]{./media/image4.png}

\section{\texorpdfstring{\textbf{v0.2 --- Triangles and
Clipping}}{v0.2 --- Triangles and Clipping}}\label{v0.2-triangles-and-clipping}

{\def\LTcaptype{none} % do not increment counter
\begin{longtable}[]{@{}
  >{\raggedright\arraybackslash}p{(\linewidth - 2\tabcolsep) * \real{0.1923}}
  >{\raggedright\arraybackslash}p{(\linewidth - 2\tabcolsep) * \real{0.8077}}@{}}
\toprule\noalign{}
\begin{minipage}[b]{\linewidth}\raggedright
\textbf{v0.2}
\end{minipage} & \begin{minipage}[b]{\linewidth}\raggedright
\textbf{Möller-Trumbore + Near/Far Clipping}
\end{minipage} \\
\begin{minipage}[b]{\linewidth}\raggedright
\textbf{Files changed}
\end{minipage} & \begin{minipage}[b]{\linewidth}\raggedright
Triangle.java, Scene.java, Vector3D.java (cross product)
\end{minipage} \\
\midrule\noalign{}
\endhead
\bottomrule\noalign{}
\endlastfoot
\end{longtable}
}

This version added the Triangle class, which is the building block for
all 3D models. Almost everything in 3D graphics is made of triangles,
any mesh, any OBJ file, is just a collection of triangles.

\subsubsection{\texorpdfstring{\textbf{Möller-Trumbore
algorithm}}{Möller-Trumbore algorithm}}\label{muxf6ller-trumbore-algorithm}

The algorithm to intersect a ray with a triangle is called
Möller-Trumbore, and it\textquotesingle s the standard approach in
raytracers because it\textquotesingle s efficient, it avoids computing
the plane equation explicitly. The idea is:

\begin{itemize}
\item
  Compute P = D × edge1 (cross product of ray direction and one edge)
\item
  The determinant = edge2 · P tells if the ray is parallel to the
  triangle
\item
  Solve for barycentric coordinates u and v --- if either is outside
  {[}0,1{]} the hit point is outside the triangle
\item
  Solve for t --- the distance along the ray to the intersection
\end{itemize}

Barycentric coordinates are important: they give you the position inside
the triangle as a weighted combination of the three vertices. This
becomes crucial later for interpolating normals and UV coordinates.

\subsubsection{\texorpdfstring{\textbf{Clipping
planes}}{Clipping planes}}\label{clipping-planes}

Clipping was also added here: a nearClip (0.1) and farClip (100.0)
define a valid range for intersections. Any hit with t \textless{}
nearClip is behind or too close to the camera, and t \textgreater{}
farClip is too far. This avoids rendering objects behind the camera and
also prevents floating-point precision issues.

\section{\texorpdfstring{\textbf{v0.3 --- OBJ File
Reader}}{v0.3 --- OBJ File Reader}}\label{v0.3-obj-file-reader}

{\def\LTcaptype{none} % do not increment counter
\begin{longtable}[]{@{}
  >{\raggedright\arraybackslash}p{(\linewidth - 2\tabcolsep) * \real{0.1923}}
  >{\raggedright\arraybackslash}p{(\linewidth - 2\tabcolsep) * \real{0.8077}}@{}}
\toprule\noalign{}
\begin{minipage}[b]{\linewidth}\raggedright
\textbf{v0.3}
\end{minipage} & \begin{minipage}[b]{\linewidth}\raggedright
\textbf{Loading 3D Models from .obj Files}
\end{minipage} \\
\begin{minipage}[b]{\linewidth}\raggedright
\textbf{Files changed}
\end{minipage} & \begin{minipage}[b]{\linewidth}\raggedright
OBJ.java
\end{minipage} \\
\midrule\noalign{}
\endhead
\bottomrule\noalign{}
\endlastfoot
\end{longtable}
}

Having triangles is one thing, but you can\textquotesingle t build a
whole scene by hand-coding every vertex. The OBJ format is a simple
text-based format that Blender and most 3D tools export. Each line
starts with a keyword: v for vertex, vn for normal, vt for texture
coordinate, f for face.

The OBJ reader parses each line and builds Triangle objects. Faces can
be quads (4 vertices) or even n-gons, so the reader does fan
triangulation, it takes the first vertex and connects it to every edge
of the polygon, splitting it into triangles.

{\def\LTcaptype{none} % do not increment counter
\begin{longtable}[]{@{}
  >{\raggedright\arraybackslash}p{(\linewidth - 0\tabcolsep) * \real{1.0000}}@{}}
\toprule\noalign{}
\begin{minipage}[b]{\linewidth}\raggedright
// Fan triangulation for quads/ngons

for (int i = 1; i \textless{} vIdx.length - 1; i++) \{

v0 = vertices{[}vIdx{[}0{]}{]};

v1 = vertices{[}vIdx{[}i{]}{]};

v2 = vertices{[}vIdx{[}i+1{]}{]};

triangles.add(new Triangle(v0, v1, v2, color));

\}
\end{minipage} \\
\midrule\noalign{}
\endhead
\bottomrule\noalign{}
\endlastfoot
\end{longtable}
}

A \textbf{scale} parameter was also added so models that are exported at
weird sizes can be \textbf{scaled} down when loading. Many OBJ files
have coordinates in the hundreds or thousands, so a scale of 0.01 brings
them to a reasonable size.

\includegraphics[width=5.10417in,height=3.82813in]{./media/image6.png}

\section{\texorpdfstring{\textbf{v0.4 --- Flat
Shading}}{v0.4 --- Flat Shading}}\label{v0.4-flat-shading}

{\def\LTcaptype{none} % do not increment counter
\begin{longtable}[]{@{}
  >{\raggedright\arraybackslash}p{(\linewidth - 2\tabcolsep) * \real{0.1923}}
  >{\raggedright\arraybackslash}p{(\linewidth - 2\tabcolsep) * \real{0.8077}}@{}}
\toprule\noalign{}
\begin{minipage}[b]{\linewidth}\raggedright
\textbf{v0.4}
\end{minipage} & \begin{minipage}[b]{\linewidth}\raggedright
\textbf{Lambert Diffuse Lighting}
\end{minipage} \\
\begin{minipage}[b]{\linewidth}\raggedright
\textbf{Files changed}
\end{minipage} & \begin{minipage}[b]{\linewidth}\raggedright
Light.java, Raytracer.java (applyShading)
\end{minipage} \\
\midrule\noalign{}
\endhead
\bottomrule\noalign{}
\endlastfoot
\end{longtable}
}

Up to this point everything was just solid flat colors and no lighting
at all. Flat shading adds a basic diffuse light calculation using
Lambert\textquotesingle s law: the brightness of a surface depends on
how directly it faces the light source.

The formula is: Diffuse = LC × OC × LI × max(0, N·L), where LC is the
light color, OC is the object color, LI is the light intensity, N is the
surface normal, and L is the direction toward the light. The dot product
N·L gives a value from -1 to 1 --- negative values mean the face is
pointing away from the light, so we clamp to 0.

A Light class was created with color and intensity. The direction of the
light determines which faces are lit and which are in shadow.

\includegraphics[width=5.17361in,height=3.88021in]{./media/image1.png}

\section{\texorpdfstring{\textbf{v0.5 --- Phong
Shading}}{v0.5 --- Phong Shading}}\label{v0.5-phong-shading}

{\def\LTcaptype{none} % do not increment counter
\begin{longtable}[]{@{}
  >{\raggedright\arraybackslash}p{(\linewidth - 2\tabcolsep) * \real{0.1923}}
  >{\raggedright\arraybackslash}p{(\linewidth - 2\tabcolsep) * \real{0.8077}}@{}}
\toprule\noalign{}
\begin{minipage}[b]{\linewidth}\raggedright
\textbf{v0.5}
\end{minipage} & \begin{minipage}[b]{\linewidth}\raggedright
\textbf{Per-vertex Normal Interpolation + Specular}
\end{minipage} \\
\begin{minipage}[b]{\linewidth}\raggedright
\textbf{Files changed}
\end{minipage} & \begin{minipage}[b]{\linewidth}\raggedright
Triangle.java, OBJ.java, Raytracer.java
\end{minipage} \\
\midrule\noalign{}
\endhead
\bottomrule\noalign{}
\endlastfoot
\end{longtable}
}

Flat shading gives a faceted look, each triangle is uniformly colored.
Phong shading smooths this out by interpolating normals across the
triangle surface. Instead of one normal per triangle, each vertex gets
its own normal (read from vn in the OBJ file), and for any point inside
the triangle the normal is interpolated using the barycentric
coordinates from Möller-Trumbore.

This makes curved surfaces look smooth even though
they\textquotesingle re made of flat triangles. The specular highlight
was also added using the reflection vector R = 2*(N·L)*N - L and the
view vector V, the specular term is (R·V)\^{}shininess. Higher shininess
makes the highlight smaller and sharper.

\includegraphics[width=4.77083in,height=3.57813in]{./media/image7.png}

\section{\texorpdfstring{\textbf{v0.6 --- Point Lights and Multiple
Lights}}{v0.6 --- Point Lights and Multiple Lights}}\label{v0.6-point-lights-and-multiple-lights}

{\def\LTcaptype{none} % do not increment counter
\begin{longtable}[]{@{}
  >{\raggedright\arraybackslash}p{(\linewidth - 2\tabcolsep) * \real{0.1923}}
  >{\raggedright\arraybackslash}p{(\linewidth - 2\tabcolsep) * \real{0.8077}}@{}}
\toprule\noalign{}
\begin{minipage}[b]{\linewidth}\raggedright
\textbf{v0.6}
\end{minipage} & \begin{minipage}[b]{\linewidth}\raggedright
\textbf{Light Hierarchy + Attenuation}
\end{minipage} \\
\begin{minipage}[b]{\linewidth}\raggedright
\textbf{Files changed}
\end{minipage} & \begin{minipage}[b]{\linewidth}\raggedright
Light.java, DirectionalLight.java, PointLight.java, Scene.java
\end{minipage} \\
\midrule\noalign{}
\endhead
\bottomrule\noalign{}
\endlastfoot
\end{longtable}
}

The original light was a simple directional light, same direction
everywhere, no falloff. Point lights are more realistic: they have a
position in 3D space and their intensity decreases with distance squared
(inverse square law, same as real physics).

The Light class became abstract with two subclasses. DirectionalLight
returns the same L direction for every point in the scene. PointLight
calculates L dynamically as normalize(lightPos - intersectionPoint) and
attenuates intensity as I / d\^{}2.

Scene.java was updated to hold a list of lights so multiple lights could
be added. The shading loop iterates over all lights and accumulates
their contributions.

\section{\texorpdfstring{\textbf{v0.7 --- Shadow
Rays}}{v0.7 --- Shadow Rays}}\label{v0.7-shadow-rays}

{\def\LTcaptype{none} % do not increment counter
\begin{longtable}[]{@{}
  >{\raggedright\arraybackslash}p{(\linewidth - 2\tabcolsep) * \real{0.1923}}
  >{\raggedright\arraybackslash}p{(\linewidth - 2\tabcolsep) * \real{0.8077}}@{}}
\toprule\noalign{}
\begin{minipage}[b]{\linewidth}\raggedright
\textbf{v0.7}
\end{minipage} & \begin{minipage}[b]{\linewidth}\raggedright
\textbf{Hard Shadows via Occlusion Testing}
\end{minipage} \\
\begin{minipage}[b]{\linewidth}\raggedright
\textbf{Files changed}
\end{minipage} & \begin{minipage}[b]{\linewidth}\raggedright
Scene.java (isOccluded), Raytracer.java
\end{minipage} \\
\midrule\noalign{}
\endhead
\bottomrule\noalign{}
\endlastfoot
\end{longtable}
}

To calculate shadows, after finding an intersection the raytracer shoots
a second ray (the shadow ray) from the hit point toward each light
source. If this ray hits any other object before reaching the light,
that point is in shadow and the light is skipped.

\subsubsection{\texorpdfstring{\textbf{Shadow
acne}}{Shadow acne}}\label{shadow-acne}

A classic problem here is shadow acne, the shadow ray can intersect the
same surface it\textquotesingle s coming from due to floating-point
precision. The fix is to offset the shadow ray origin slightly along the
normal by a small bias (0.001), so it doesn\textquotesingle t
immediately re-intersect its own surface.

For point lights, the shadow ray also needs a maximum distance, only
objects between the surface and the light should cast shadows. Objects
behind the light shouldn\textquotesingle t matter.

\section{\texorpdfstring{\textbf{v0.75 --- Textures from MTL
Files}}{v0.75 --- Textures from MTL Files}}\label{v0.75-textures-from-mtl-files}

{\def\LTcaptype{none} % do not increment counter
\begin{longtable}[]{@{}
  >{\raggedright\arraybackslash}p{(\linewidth - 2\tabcolsep) * \real{0.1923}}
  >{\raggedright\arraybackslash}p{(\linewidth - 2\tabcolsep) * \real{0.8077}}@{}}
\toprule\noalign{}
\begin{minipage}[b]{\linewidth}\raggedright
\textbf{v0.75}
\end{minipage} & \begin{minipage}[b]{\linewidth}\raggedright
\textbf{UV Mapping + Diffuse Texture Sampling}
\end{minipage} \\
\begin{minipage}[b]{\linewidth}\raggedright
\textbf{Files changed}
\end{minipage} & \begin{minipage}[b]{\linewidth}\raggedright
Material.java, MTL.java, Triangle.java, OBJ.java, Intersection.java,
Raytracer.java
\end{minipage} \\
\midrule\noalign{}
\endhead
\bottomrule\noalign{}
\endlastfoot
\end{longtable}
}

This was a big architectural step. The OBJ format has a companion .mtl
file that defines materials --- including paths to texture images. The
texture is a PNG or JPG that gets mapped onto the surface using UV
coordinates.

UV coordinates (vt lines in the OBJ) are 2D coordinates in {[}0,1{]}
range that map each vertex to a position on the texture image. Just like
normals, they get interpolated with barycentric coordinates for every
pixel inside a triangle. Then sampleDiffuse(u, v) does a simple pixel
lookup on the image.

Intersection.java was updated to carry the interpolated u and v values.
Material.java handles loading the image with ImageIO and sampling it.
MTL.java parses the .mtl file and creates Material objects for each
material name found.

\includegraphics[width=5.29861in,height=3.97396in]{./media/image2.png}

\section{\texorpdfstring{\textbf{v0.76 ---
BVH}}{v0.76 --- BVH}}\label{v0.76-bvh}

{\def\LTcaptype{none} % do not increment counter
\begin{longtable}[]{@{}
  >{\raggedright\arraybackslash}p{(\linewidth - 2\tabcolsep) * \real{0.1923}}
  >{\raggedright\arraybackslash}p{(\linewidth - 2\tabcolsep) * \real{0.8077}}@{}}
\toprule\noalign{}
\begin{minipage}[b]{\linewidth}\raggedright
\textbf{v0.76}
\end{minipage} & \begin{minipage}[b]{\linewidth}\raggedright
\textbf{Bounding Volume Hierarchy}
\end{minipage} \\
\begin{minipage}[b]{\linewidth}\raggedright
\textbf{Files changed}
\end{minipage} & \begin{minipage}[b]{\linewidth}\raggedright
AABB.java, BVHNode.java, Scene.java
\end{minipage} \\
\midrule\noalign{}
\endhead
\bottomrule\noalign{}
\endlastfoot
\end{longtable}
}

Without any acceleration structure, every ray has to be tested against
every single triangle in the scene. A scene with 100,000 triangles at
1080p resolution means about 2 trillion intersection tests per frame,
and it is completely impractical.

The BVH (Bounding Volume Hierarchy) solves this by organizing objects
into a tree. Each node has an AABB (axis-aligned bounding box) that
contains all its children. If a ray misses the bounding box, the entire
subtree is skipped.

\subsubsection{\texorpdfstring{\textbf{Construction}}{Construction}}\label{construction}

The tree is built by: computing the AABB of all objects, finding the
longest axis, sorting objects by their centroid on that axis, splitting
at the median, and recursing on each half. Leaves contain at most 8
objects. This is O(n log n) and happens once before rendering.

\subsubsection{\texorpdfstring{\textbf{AABB intersection (the slab
method)}}{AABB intersection (the slab method)}}\label{aabb-intersection-the-slab-method}

For each axis, compute when the ray enters (t1) and exits (t2) the pair
of parallel planes. Track the latest entry and earliest exit across all
three axes. If the latest entry is before the earliest exit, the ray hit
the box. Pure arithmetic, no square roots.

\section{\texorpdfstring{\textbf{v0.77 --- Normal Maps (Bump
Maps)}}{v0.77 --- Normal Maps (Bump Maps)}}\label{v0.77-normal-maps-bump-maps}

{\def\LTcaptype{none} % do not increment counter
\begin{longtable}[]{@{}
  >{\raggedright\arraybackslash}p{(\linewidth - 2\tabcolsep) * \real{0.1923}}
  >{\raggedright\arraybackslash}p{(\linewidth - 2\tabcolsep) * \real{0.8077}}@{}}
\toprule\noalign{}
\begin{minipage}[b]{\linewidth}\raggedright
\textbf{v0.77}
\end{minipage} & \begin{minipage}[b]{\linewidth}\raggedright
\textbf{TBN Matrix + Tangent-Space Normal Perturbation}
\end{minipage} \\
\begin{minipage}[b]{\linewidth}\raggedright
\textbf{Files changed}
\end{minipage} & \begin{minipage}[b]{\linewidth}\raggedright
Material.java, MTL.java, Triangle.java, Raytracer.java
\end{minipage} \\
\midrule\noalign{}
\endhead
\bottomrule\noalign{}
\endlastfoot
\end{longtable}
}

Normal maps allow a low-polygon mesh to look like it has much more
geometric detail --- things like brick grout, surface scratches, or skin
pores --- without actually adding geometry. The trick is to store
per-pixel normal directions in a texture and use them to perturb the
shading normal.

\subsubsection{\texorpdfstring{\textbf{The TBN
matrix}}{The TBN matrix}}\label{the-tbn-matrix}

The normal map stores vectors in tangent space --- a coordinate system
local to each triangle where Z is the surface normal, X is the tangent
(along U in UV space), and Y is the bitangent (along V in UV space). To
use these in world space, a 3x3 TBN matrix is built from the tangent
(T), bitangent (B), and normal (N) vectors.

T and B are calculated from the UV differences and edge vectors of the
triangle. This math is done once per triangle in the constructor:

{\def\LTcaptype{none} % do not increment counter
\begin{longtable}[]{@{}
  >{\raggedright\arraybackslash}p{(\linewidth - 0\tabcolsep) * \real{1.0000}}@{}}
\toprule\noalign{}
\begin{minipage}[b]{\linewidth}\raggedright
du1 = uv1{[}0{]}-uv0{[}0{]}; dv1 = uv1{[}1{]}-uv0{[}1{]};

du2 = uv2{[}0{]}-uv0{[}0{]}; dv2 = uv2{[}1{]}-uv0{[}1{]};

f = 1.0 / (du1*dv2 - du2*dv1);

tangent = normalize(f * (dv2*edge1 - dv1*edge2));

bitangent = normalize(f * (-du2*edge1 + du1*edge2));
\end{minipage} \\
\midrule\noalign{}
\endhead
\bottomrule\noalign{}
\endlastfoot
\end{longtable}
}

At shading time, the RGB values from the normal map are decoded from
{[}0,255{]} to {[}-1,1{]}, then transformed to world space: N\_world =
tsN.x*T + tsN.y*B + tsN.z*N. Before this, T is re-orthogonalized against
the Phong-interpolated normal using Gram-Schmidt to keep everything
consistent.

\includegraphics[width=6.5in,height=3.43056in]{./media/image12.png}

\section{\texorpdfstring{\textbf{v0.8 ---
Blinn-Phong}}{v0.8 --- Blinn-Phong}}\label{v0.8-blinn-phong}

{\def\LTcaptype{none} % do not increment counter
\begin{longtable}[]{@{}
  >{\raggedright\arraybackslash}p{(\linewidth - 2\tabcolsep) * \real{0.1923}}
  >{\raggedright\arraybackslash}p{(\linewidth - 2\tabcolsep) * \real{0.8077}}@{}}
\toprule\noalign{}
\begin{minipage}[b]{\linewidth}\raggedright
\textbf{v0.8}
\end{minipage} & \begin{minipage}[b]{\linewidth}\raggedright
\textbf{Hybrid Material Architecture}
\end{minipage} \\
\begin{minipage}[b]{\linewidth}\raggedright
\textbf{Files changed}
\end{minipage} & \begin{minipage}[b]{\linewidth}\raggedright
BlinnPhongMaterial.java, TypicalMaterials.java, Triangle.java,
Scene.java, Raytracer.java
\end{minipage} \\
\midrule\noalign{}
\endhead
\bottomrule\noalign{}
\endlastfoot
\end{longtable}
}

Up to this point, light behavior (how shiny, how diffuse) was hardcoded
in the shader. This version introduced a proper material system that
separates two distinct concerns:

\begin{itemize}
\item
  Material.java --- only holds textures from the MTL file (diffuse
  texture, normal map)
\item
  BlinnPhongMaterial.java --- only holds light behavior coefficients
  (ka, kd, ks, shininess, transparency, reflectivity, ior)
\end{itemize}

A Triangle can have both: its MTL Material provides the color and normal
map, and a BlinnPhongMaterial is assigned when adding the mesh to the
scene to control how it responds to light. This means you can load the
same OBJ with different material behaviors without modifying the
geometry.

\subsubsection{\texorpdfstring{\textbf{Blinn vs Phong
specular}}{Blinn vs Phong specular}}\label{blinn-vs-phong-specular}

The specular calculation was upgraded from classic Phong to Blinn-Phong.
Instead of computing the reflect vector R and doing dot(R, V), it
computes a half-vector H = normalize(L + V) and does dot(N, H). This is
slightly cheaper and avoids artifacts at grazing angles. The
TypicalMaterials class provides preset materials like GOLD, SILVER,
IRON, GLASS, MARBLE, DIAMOND with realistic coefficient values.

\includegraphics[width=4.89063in,height=2.75098in]{./media/image10.png}

\section{\texorpdfstring{\textbf{v0.9 --- Recursive
Reflections}}{v0.9 --- Recursive Reflections}}\label{v0.9-recursive-reflections}

{\def\LTcaptype{none} % do not increment counter
\begin{longtable}[]{@{}
  >{\raggedright\arraybackslash}p{(\linewidth - 2\tabcolsep) * \real{0.1923}}
  >{\raggedright\arraybackslash}p{(\linewidth - 2\tabcolsep) * \real{0.8077}}@{}}
\toprule\noalign{}
\begin{minipage}[b]{\linewidth}\raggedright
\textbf{v0.9}
\end{minipage} & \begin{minipage}[b]{\linewidth}\raggedright
\textbf{Mirror Reflections with Bounce Limit}
\end{minipage} \\
\begin{minipage}[b]{\linewidth}\raggedright
\textbf{Files changed}
\end{minipage} & \begin{minipage}[b]{\linewidth}\raggedright
Raytracer.java
\end{minipage} \\
\midrule\noalign{}
\endhead
\bottomrule\noalign{}
\endlastfoot
\end{longtable}
}

Reflections are what make raytracers look so much better than
rasterizers. When a ray hits a reflective surface, instead of stopping,
it spawns a new ray in the reflected direction and traces that
recursively.

The reflected direction is: R = I - 2*(I·N)*N, where I is the incident
ray direction. A MAX\_BOUNCES constant (default 2-3) limits the
recursion depth to avoid infinite bouncing between two mirrors.

The final color blends the local shading with the reflected color:
finalColor = (1-kr)*localColor + kr*reflectedColor. With kr=1.0 and no
diffuse, you get a perfect mirror.

\includegraphics[width=4.69271in,height=2.63965in]{./media/image8.png}\includegraphics[width=4.66146in,height=2.62207in]{./media/image5.png}

\section{\texorpdfstring{\textbf{v1.0 --- Refraction with
Snell\textquotesingle s Law and
Fresnel}}{v1.0 --- Refraction with Snell\textquotesingle s Law and Fresnel}}\label{v1.0-refraction-with-snells-law-and-fresnel}

{\def\LTcaptype{none} % do not increment counter
\begin{longtable}[]{@{}
  >{\raggedright\arraybackslash}p{(\linewidth - 2\tabcolsep) * \real{0.1923}}
  >{\raggedright\arraybackslash}p{(\linewidth - 2\tabcolsep) * \real{0.8077}}@{}}
\toprule\noalign{}
\begin{minipage}[b]{\linewidth}\raggedright
\textbf{v1.0}
\end{minipage} & \begin{minipage}[b]{\linewidth}\raggedright
\textbf{Transparent Materials + IOR + Fresnel}
\end{minipage} \\
\begin{minipage}[b]{\linewidth}\raggedright
\textbf{Files changed}
\end{minipage} & \begin{minipage}[b]{\linewidth}\raggedright
BlinnPhongMaterial.java, TypicalMaterials.java, Raytracer.java
\end{minipage} \\
\midrule\noalign{}
\endhead
\bottomrule\noalign{}
\endlastfoot
\end{longtable}
}

Refraction is what makes glass, water, and gems look like glass, water,
and gems. When a ray enters a transparent material, it bends according
to Snell\textquotesingle s Law: n1*sin(θ1) = n2*sin(θ2), where n1 and n2
are the refractive indices (IOR) of the two materials.

\subsubsection{\texorpdfstring{\textbf{IOR --- Index of
Refraction}}{IOR --- Index of Refraction}}\label{ior-index-of-refraction}

The IOR was added as a seventh parameter to BlinnPhongMaterial.
Real-world values were used in TypicalMaterials: GLASS=1.50, WATER=1.33,
DIAMOND=2.42, EMERALD=1.57. Higher IOR means more bending and more
distortion.

\subsubsection{\texorpdfstring{\textbf{Snell\textquotesingle s Law in
vector
form}}{Snell\textquotesingle s Law in vector form}}\label{snells-law-in-vector-form}

{\def\LTcaptype{none} % do not increment counter
\begin{longtable}[]{@{}
  >{\raggedright\arraybackslash}p{(\linewidth - 0\tabcolsep) * \real{1.0000}}@{}}
\toprule\noalign{}
\begin{minipage}[b]{\linewidth}\raggedright
// cosI = -dot(N, I) (angle of incidence)

// ratio = n1/n2

// sinT2 = ratio\^{}2 * (1 - cosI\^{}2)

// if sinT2 \textgreater{} 1 → Total Internal Reflection (no refracted
ray)

cosT = sqrt(1 - sinT2);

refractDir = ratio*I + (ratio*cosI - cosT)*N;
\end{minipage} \\
\midrule\noalign{}
\endhead
\bottomrule\noalign{}
\endlastfoot
\end{longtable}
}

\subsubsection{\texorpdfstring{\textbf{Fresnel effect (Schlick
approximation)}}{Fresnel effect (Schlick approximation)}}\label{fresnel-effect-schlick-approximation}

Even transparent materials reflect more at grazing angles --- this is
the Fresnel effect, and it\textquotesingle s why glass looks transparent
from straight on but mirror-like from the side. The Schlick
approximation computes this efficiently:

{\def\LTcaptype{none} % do not increment counter
\begin{longtable}[]{@{}
  >{\raggedright\arraybackslash}p{(\linewidth - 0\tabcolsep) * \real{1.0000}}@{}}
\toprule\noalign{}
\begin{minipage}[b]{\linewidth}\raggedright
r0 = ((n1-n2)/(n1+n2))\^{}2;

fresnelKr = r0 + (1 - r0) * (1 - cosI)\^{}5;
\end{minipage} \\
\midrule\noalign{}
\endhead
\bottomrule\noalign{}
\endlastfoot
\end{longtable}
}

The final color blend is: finalColor = localColor*(1-kr-kt) +
reflectedColor*kr + refractedColor*kt. The three contributions are
energy-conserving --- they sum to at most 1.

Total Internal Reflection is also handled: when a ray inside a dense
material hits the boundary at too steep an angle, sin\^{}2(T)
\textgreater{} 1 and there is no refracted ray --- all light reflects
back inside. This is what causes the characteristic sparkling in
diamonds and gems.

\includegraphics[width=5.30729in,height=3.00236in]{./media/image11.png}

\section{\texorpdfstring{\textbf{Final
Architecture}}{Final Architecture}}\label{final-architecture}

After all iterations, the codebase is organized as follows:

\subsubsection{\texorpdfstring{\textbf{Core
math}}{Core math}}\label{core-math}

\begin{itemize}
\item
  Vector3D.java --- all vector math (add, sub, dot, cross, normalize,
  magnitude)
\item
  Ray.java --- origin + direction, with pre-computed inverse direction
  for AABB
\end{itemize}

\subsubsection{\texorpdfstring{\textbf{Scene and
acceleration}}{Scene and acceleration}}\label{scene-and-acceleration}

\begin{itemize}
\item
  Scene.java --- holds objects, lights, camera, BVH. Main raycast and
  shadow test entry points
\item
  AABB.java --- axis-aligned bounding box with slab-method ray
  intersection
\item
  BVHNode.java --- binary tree of AABBs built from scene objects
\end{itemize}

\subsubsection{\texorpdfstring{\textbf{Geometry}}{Geometry}}\label{geometry}

\begin{itemize}
\item
  Object3D.java --- abstract base class
\item
  Triangle.java --- Möller-Trumbore, Phong normals, UV interpolation,
  TBN, BlinnPhongMaterial
\item
  Sphere.java --- analytic ray-sphere intersection
\item
  Plane.java --- infinite plane for floors and walls
\end{itemize}

\subsubsection{\texorpdfstring{\textbf{Materials and
textures}}{Materials and textures}}\label{materials-and-textures}

\begin{itemize}
\item
  Material.java --- diffuse texture + normal map loading and sampling
\item
  MTL.java --- .mtl file parser (map\_Kd and map\_Bump)
\item
  BlinnPhongMaterial.java --- ka, kd, ks, shininess, transparency,
  reflectivity, IOR
\item
  TypicalMaterials.java --- presets: GOLD, SILVER, GLASS, DIAMOND,
  STONE, etc.
\end{itemize}

\subsubsection{\texorpdfstring{\textbf{Lighting}}{Lighting}}\label{lighting}

\begin{itemize}
\item
  Light.java --- abstract base
\item
  DirectionalLight.java --- constant direction, no attenuation
\item
  PointLight.java --- position-based, intensity / d\^{}2 attenuation
\end{itemize}

\subsubsection{\texorpdfstring{\textbf{Rendering}}{Rendering}}\label{rendering}

\begin{itemize}
\item
  Camera.java --- ray generation from pixel coordinates
\item
  Raytracer.java --- main render loop, traceRay() recursive function,
  applyShading() with full Blinn-Phong + normal maps + reflection +
  refraction + Fresnel
\end{itemize}

\section{FINAL RENDERS AND
EXPLANATION}\label{final-renders-and-explanation}

\textbf{RENDER 1 -- The Temple and The Legend}

With the first render, I wanted to do something I like, and what this
render aims to tell is the phrase "all games have a story, but only one
is a legend," and that\textquotesingle s the story: the legend turned
out to be real, and the sword rests on its pedestal, waiting to be
discovered by the hero.

\includegraphics[width=8.21574in,height=4.35938in]{./media/image9.png}

\textbf{RENDER 2 -- A Good Beginning\ldots{}}

We\textquotesingle ve all started something for the first time. I
suppose almost everyone here knows Minecraft, and this render tries to
convey the feeling of starting something for the first time and having
everything go well (Steve just started and already has diamonds).

Which doesn\textquotesingle t mean it will be perfect later on
(represented by Herobrine in the mirror\textquotesingle s reflection).

\includegraphics[width=7.95022in,height=4.19271in]{./media/image3.png}

\textbf{RENDER 3 -- Knight\textquotesingle s Dream}

A knight fell defeated in his greatest battle, but all is not lost. The
knight entered a dream where an angel appeared to him. The knight asked
for help and strength to continue. This is the scene depicted.

The aim is to tell, through more fantastical objects and more saturated
colors, the effect of a dream---a dream that will bring something good.

\includegraphics[width=8.13021in,height=4.2866in]{./media/image13.png}

\section{EXPLANATION OF GREATEST
CHALLENGE}\label{explanation-of-greatest-challenge}

My biggest challenge was the phong shading, it looked weird and also
like flat shading.

The problem was in the Möller--Trumbore algorithm, more specifically in
the barycentric coordinates, I had inverted the V and U, so it scrambles
the normals between vertices, producing incorrect interpolation that
visually resembles flat shading or gives strange results.

The fix was given by a friend, that has the same problem, and say me
that to fix that, i only need to invert the U and V in the formula

\section{SELF EVALUATION}\label{self-evaluation}

\subsubsection{Ray Tracer v1.0 (10\%)}\label{ray-tracer-v1.0-10}

\begin{itemize}
\item
  Shadows (5\%)
\item
  Blinn-Phong (specular \& shininess) (5\%)
\end{itemize}

\ul{In this part I would give me the 10\% because my shadows and
Blinn-Phong are working correctly.}

\subsubsection{Reflection Works (25\%)}\label{reflection-works-25}

\begin{itemize}
\item
  Reflected objects are shown correctly (15\%)
\item
  It shows at least 2 bounces (10\%)
\end{itemize}

\ul{I also would give me the 25\%, also the reflections work as
expected}

\subsubsection{Refraction Works (25\%)}\label{refraction-works-25}

\begin{itemize}
\item
  Refracted objects are shown correctly (15\%)
\item
  Light changes direction correctly (10\%)
\end{itemize}

\ul{I also would give me the 25\%, also the Refraction work}

\subsubsection{Coding Principles (15\%)}\label{coding-principles-15}

\begin{itemize}
\item
  Does it have unnecessary code? (1\%)
\item
  Does it properly use OOP? (2\%)
\item
  Is the code reusable? (1\%)
\item
  Is the code flexible? (1\%)
\item
  Does it have bugs? (2\%)
\item
  Is the code scalable? (1\%)
\item
  Does it have comments? (1\%)
\item
  Is the code a huge mess or neat? (1\%)
\end{itemize}

\ul{Here i would give me a 14\% or 13\%, because i think that my code is
a little bit a mess, especially in the part of the creations of the
scenes, i have a lot of code (i think that is organized, but also i
think it can be done in a cleaner way)}

\subsubsection{LaTeX Report (5\%)}\label{latex-report-5}

A full report in LaTeX explaining how you created the ray tracer and
showing at least 5 images of how your ray tracer evolved. Include the
explanation for each of the final images.

Evaluation:

\begin{itemize}
\item
  Description of the process of the ray tracer (1.5\%)
\item
  Explanation of each of the final images (1.5\%)
\item
  Explanation of greatest challenges (1\%)
\item
  Self-evaluation of all points (1\%)
\end{itemize}

\ul{In this part i would give me a 5\%, because i do all the points,
including the process of the ray tracer, i explain each of the final
images, i explain my greatest challenges and i auto-evaluate my self}

\subsubsection{3 Complex Final Renders
(10\%)}\label{complex-final-renders-10}

\begin{itemize}
\item
  Each render has a story (3\%)
\item
  Each render measures 4096 × 2160 and is PNG (2\%)
\item
  It contains something more than spheres or teapots (2\%)
\item
  It is aesthetic (3\%)
\end{itemize}

\ul{here also i give me the 10\%, my renders tell a story and also are
4K, they also have more than spheres and teapots and in my opinion its
aesthetic.}

\textbf{CHALLENGE MODE}

\ul{I think I can get 1.5\% from the challenge mode of ``Implement
Rotation and Scale of Objects'', because I implement the scale of the
objects.}

\ul{I also would give me 5\% of the extra points, because i also
implement the Bounding Volume Hierarchy}

\ul{And i also implement the UV reader + texture and normals maps, so i
would give me that 8\% extra}

\end{document}
